I agree with the bug report that several “standard facts” should be written out explicitly.  
In particular, I will **not** assume in advance that the support points \(x_i=s_i+t_i\) are distinct; that will follow from the proof.

Let
\[
R:=P+Q,\qquad E_i:=[s_i,s_{i+1}],\qquad a_i:=s_{i+1}-s_i,
\]
with indices taken modulo \(n\).

Since \(P=\operatorname{conv}(S)\) and \(Q=\operatorname{conv}(T)\),
\[
R=P+Q=\operatorname{conv}(S+T),
\]
so \(R\) is a convex polygon.

---

## Preliminary lemmas

### Lemma 1: exposed faces of a convex hull of finitely many points
If \(A\subset \mathbb{R}^2\) is finite and \(K=\operatorname{conv}(A)\), define
\[
M_A(u):=\left\{a\in A:\ \langle u,a\rangle=\max_{x\in A}\langle u,x\rangle\right\}.
\]
Then
\[
F_K(u)=\operatorname{conv}\bigl(M_A(u)\bigr).
\]

#### Proof
Let
\[
m:=\max_{x\in A}\langle u,x\rangle.
\]
If \(x=\sum_{a\in A}\alpha_a a\in K\) with \(\alpha_a\ge 0\) and \(\sum_a\alpha_a=1\), then
\[
\langle u,x\rangle=\sum_{a\in A}\alpha_a\langle u,a\rangle\le \sum_{a\in A}\alpha_a m=m.
\]
So \(h_K(u)=m\).

Moreover, equality \(\langle u,x\rangle=m\) holds iff every \(a\) with \(\alpha_a>0\) satisfies \(\langle u,a\rangle=m\), i.e. iff \(x\in \operatorname{conv}(M_A(u))\). Hence
\[
F_K(u)=\operatorname{conv}(M_A(u)).
\]
\(\square\)

Two consequences will be used repeatedly:

- if one point of \(A\) strictly beats all others for the functional \(\langle u,\cdot\rangle\), then it is the **unique** maximizer on \(K=\operatorname{conv}(A)\);
- if \(a\in A\) merely maximizes on \(A\), then \(a\in F_K(u)\).

---

### Lemma 2: exposed faces of a Minkowski sum
For nonempty compact convex sets \(A,B\subset \mathbb{R}^2\),
\[
F_{A+B}(u)=F_A(u)+F_B(u).
\]

#### Proof
First,
\[
h_{A+B}(u)=\max_{a\in A,\ b\in B}\langle u,a+b\rangle
=\max_{a,b}\bigl(\langle u,a\rangle+\langle u,b\rangle\bigr)
=h_A(u)+h_B(u).
\]
If \(a\in F_A(u)\) and \(b\in F_B(u)\), then
\[
\langle u,a+b\rangle=h_A(u)+h_B(u)=h_{A+B}(u),
\]
so \(a+b\in F_{A+B}(u)\). Thus
\[
F_A(u)+F_B(u)\subseteq F_{A+B}(u).
\]

Conversely, if \(z\in F_{A+B}(u)\), write \(z=a+b\) with \(a\in A\), \(b\in B\). Then
\[
h_A(u)+h_B(u)=h_{A+B}(u)=\langle u,z\rangle
=\langle u,a\rangle+\langle u,b\rangle\le h_A(u)+h_B(u),
\]
so equality holds in both inequalities; hence
\[
a\in F_A(u),\qquad b\in F_B(u).
\]
Therefore \(z\in F_A(u)+F_B(u)\). \(\square\)

---

### Lemma 3: compact convex subsets of a line
Let \(a\neq 0\), and let \(G\) be a nonempty compact convex subset of a line parallel to \(a\).  
If \(p,q\in G\) are respectively the minimum and maximum points of the linear functional \(\langle a,\cdot\rangle\) on \(G\), then
\[
G=[p,q].
\]

#### Proof
Write the line as \(L=x_0+\mathbb{R}a\). Since \(G\subset L\) is compact and convex,
\[
G=\{x_0+\tau a:\ \tau\in [\alpha,\beta]\}
\]
for some \(\alpha\le \beta\).

Now
\[
\langle a, x_0+\tau a\rangle=\langle a,x_0\rangle+\tau\|a\|^2,
\]
which is strictly increasing in \(\tau\). Hence the minimum is attained at \(\tau=\alpha\) and the maximum at \(\tau=\beta\). Therefore
\[
p=x_0+\alpha a,\qquad q=x_0+\beta a,
\]
and so
\[
G=[p,q].
\]
\(\square\)

---

## Part (2)

Fix \(i\), and choose an outer normal \(u_i\) to the edge
\[
E_i=[s_i,s_{i+1}].
\]
Then
\[
F_P(u_i)=E_i,
\]
so
\[
\langle u_i,s_i\rangle=\langle u_i,s_{i+1}\rangle=h_P(u_i),
\]
and for every \(j\neq i,i+1\),
\[
\langle u_i,s_j\rangle<h_P(u_i).
\]

Let
\[
G_i:=F_Q(u_i).
\]

We prove that \(G_i=[t_i,t_{i+1}]\).

---

### Step 1: perturbing the edge normal enters the adjacent cones
For \(j\neq i,i+1\), define
\[
\delta_j:=\langle u_i,s_i-s_j\rangle>0,
\qquad
\delta:=\min_{j\neq i,i+1}\delta_j>0.
\]
Also set
\[
M:=\max_{j\neq i,i+1}\max\Bigl(
|\langle a_i,s_i-s_j\rangle|,
|\langle a_i,s_{i+1}-s_j\rangle|
\Bigr).
\]
Choose
\[
0<\varepsilon<\frac{\delta}{M+1}.
\]

Then for \(j\neq i,i+1\),
\[
\langle u_i-\varepsilon a_i,\ s_i-s_j\rangle
=\delta_j-\varepsilon\langle a_i,s_i-s_j\rangle
\ge \delta-\varepsilon M>0,
\]
and
\[
\langle u_i-\varepsilon a_i,\ s_i-s_{i+1}\rangle
=-\varepsilon\langle a_i,s_i-s_{i+1}\rangle
=\varepsilon\|a_i\|^2>0.
\]
So \(s_i\) strictly beats every other vertex of \(P\) for the functional \(\langle u_i-\varepsilon a_i,\cdot\rangle\). By Lemma 1,
\[
F_P(u_i-\varepsilon a_i)=\{s_i\}.
\]
Thus
\[
u_i-\varepsilon a_i\in C_i.
\]

Similarly, for \(j\neq i,i+1\),
\[
\langle u_i+\varepsilon a_i,\ s_{i+1}-s_j\rangle
=\langle u_i,s_{i+1}-s_j\rangle+\varepsilon\langle a_i,s_{i+1}-s_j\rangle
\ge \delta-\varepsilon M>0,
\]
and
\[
\langle u_i+\varepsilon a_i,\ s_{i+1}-s_i\rangle
=\varepsilon\|a_i\|^2>0.
\]
Hence, again by Lemma 1,
\[
F_P(u_i+\varepsilon a_i)=\{s_{i+1}\},
\]
so
\[
u_i+\varepsilon a_i\in C_{i+1}.
\]

---

### Step 2: \(t_i,t_{i+1}\in G_i\)
Take a sequence \(\varepsilon_m\downarrow 0\) such that for every \(m\),
\[
u_i-\varepsilon_m a_i\in C_i,\qquad u_i+\varepsilon_m a_i\in C_{i+1}.
\]
By the hypothesis from Subproblems \(1\)–\(2\), \(t_i\) is the unique maximizer on \(T\) for every direction in \(C_i\). Hence for every \(t\in T\),
\[
\langle u_i-\varepsilon_m a_i,\ t_i-t\rangle\ge 0.
\]
Letting \(m\to\infty\) gives
\[
\langle u_i,\ t_i-t\rangle\ge 0\qquad(\forall t\in T).
\]
So \(t_i\in M_T(u_i)\), and therefore by Lemma 1,
\[
t_i\in F_Q(u_i)=G_i.
\]

Exactly the same argument with \(C_{i+1}\) shows
\[
t_{i+1}\in G_i.
\]

---

### Step 3: \(G_i\) is exactly \([t_i,t_{i+1}]\)
Since \(G_i=F_Q(u_i)\), it lies in the support line
\[
L_i:=\{q\in\mathbb{R}^2:\ \langle u_i,q\rangle=h_Q(u_i)\}.
\]
Because \(u_i\) is orthogonal to the edge direction \(a_i\), the line \(L_i\) is parallel to \(a_i\).

Now fix \(\varepsilon>0\) as in Step 1. For any \(q\in G_i\), we have \(\langle u_i,q\rangle=h_Q(u_i)\), so
\[
\langle u_i-\varepsilon a_i,q\rangle
=h_Q(u_i)-\varepsilon\langle a_i,q\rangle.
\]
Thus, **restricted to \(G_i\)**, maximizing \(\langle u_i-\varepsilon a_i,\cdot\rangle\) is exactly the same as minimizing \(\langle a_i,\cdot\rangle\).

But \(u_i-\varepsilon a_i\in C_i\), so by the setup and Lemma 1,
\[
F_Q(u_i-\varepsilon a_i)=\{t_i\}.
\]
Since \(t_i\in G_i\), it follows that \(t_i\) is the unique minimizer of \(\langle a_i,\cdot\rangle\) on \(G_i\).

Likewise, because \(u_i+\varepsilon a_i\in C_{i+1}\),
\[
F_Q(u_i+\varepsilon a_i)=\{t_{i+1}\},
\]
and on \(G_i\),
\[
\langle u_i+\varepsilon a_i,q\rangle
=h_Q(u_i)+\varepsilon\langle a_i,q\rangle,
\]
so \(t_{i+1}\) is the unique maximizer of \(\langle a_i,\cdot\rangle\) on \(G_i\).

Now \(G_i\) is a compact convex subset of the line \(L_i\), which is parallel to \(a_i\). By Lemma 3,
\[
G_i=[t_i,t_{i+1}].
\]

Since multiplying a normal by a positive scalar does not change the exposed face, the same conclusion holds for **any** outer normal to \(E_i\).

---

### Step 4: direction of \(t_{i+1}-t_i\)
Because \(G_i=[t_i,t_{i+1}]\subset L_i\) and \(L_i\parallel a_i\), there exists \(\lambda_i\in\mathbb{R}\) such that
\[
t_{i+1}-t_i=\lambda_i a_i=\lambda_i(s_{i+1}-s_i).
\]
Also \(t_i\) is the minimum and \(t_{i+1}\) the maximum of \(\langle a_i,\cdot\rangle\) on \(G_i\), so
\[
\langle a_i,t_{i+1}-t_i\rangle\ge 0.
\]
Substituting \(t_{i+1}-t_i=\lambda_i a_i\) yields
\[
\lambda_i\|a_i\|^2\ge 0.
\]
Since \(a_i\neq 0\), we get
\[
\lambda_i\ge 0.
\]

This proves Part \((2)\).

---

## Part (1)

Define
\[
x_i:=s_i+t_i.
\]

---

### Step 1: each \(x_i\) is a vertex of \(R=P+Q\)
If \(u\in C_i\), then by definition of \(C_i\),
\[
F_P(u)=\{s_i\}.
\]
By the setup from Subproblems \(1\)–\(2\), \(t_i\) is the unique maximizer on \(T\) for every \(u\in C_i\), so by Lemma 1,
\[
F_Q(u)=\{t_i\}.
\]
Hence by Lemma 2,
\[
F_R(u)=F_P(u)+F_Q(u)=\{s_i+t_i\}=\{x_i\}.
\]
Therefore \(x_i\) is a vertex of \(R\).

---

### Step 2: for each \(i\), the corresponding edge of \(R\) is \([x_i,x_{i+1}]\)
From Part \((2)\),
\[
F_Q(u_i)=[t_i,t_{i+1}],
\]
so by Lemma 2,
\[
F_R(u_i)=E_i+[t_i,t_{i+1}].
\]

Write
\[
t_{i+1}-t_i=\lambda_i a_i,\qquad \lambda_i\ge 0.
\]
Then
\[
E_i=\{s_i+\alpha a_i:\ 0\le \alpha\le 1\},
\]
and
\[
[t_i,t_{i+1}]
=\{t_i+\beta\lambda_i a_i:\ 0\le \beta\le 1\}.
\]
Therefore
\[
E_i+[t_i,t_{i+1}]
=\{x_i+(\alpha+\lambda_i\beta)a_i:\ \alpha,\beta\in[0,1]\}.
\]

Now
\[
\{\alpha+\lambda_i\beta:\ \alpha,\beta\in[0,1]\}=[0,1+\lambda_i].
\]
Indeed:

- the inclusion \(\subseteq [0,1+\lambda_i]\) is immediate;
- conversely, if \(r\in[0,1+\lambda_i]\), then:
  - if \(0\le r\le 1\), choose \((\alpha,\beta)=(r,0)\);
  - if \(1\le r\le 1+\lambda_i\) and \(\lambda_i>0\), choose
    \[
    (\alpha,\beta)=\left(1,\frac{r-1}{\lambda_i}\right);
    \]
  - if \(\lambda_i=0\), then \([0,1+\lambda_i]=[0,1]\), already covered by the first case.

So
\[
F_R(u_i)=\{x_i+ra_i:\ 0\le r\le 1+\lambda_i\}
=[x_i,x_i+(1+\lambda_i)a_i].
\]
But
\[
x_{i+1}
=s_{i+1}+t_{i+1}
=(s_i+a_i)+(t_i+\lambda_i a_i)
=x_i+(1+\lambda_i)a_i.
\]
Hence
\[
F_R(u_i)=[x_i,x_{i+1}].
\]

Since \(1+\lambda_i>0\) and \(a_i\neq 0\), this is a nondegenerate edge of \(R\).

---

### Step 3: these are all the edges of \(R\)
Let \(E\) be any edge of \(R\), and let \(u\) be an outer normal to \(E\). Then
\[
E=F_R(u).
\]

Because \(P\) is a convex polygon, \(F_P(u)\) is either a vertex \(\{s_i\}\) or an edge \(E_i\).

- If \(F_P(u)=\{s_i\}\), then \(u\in C_i\). By the setup, \(F_Q(u)=\{t_i\}\), so Lemma 2 gives
  \[
  F_R(u)=\{x_i\},
  \]
  a point, not an edge. Contradiction.

- Therefore \(F_P(u)=E_i\) for some \(i\). Then \(u\) is an outer normal to \(E_i\), so by Part \((2)\),
  \[
  F_Q(u)=[t_i,t_{i+1}],
  \]
  and therefore
  \[
  E=F_R(u)=E_i+[t_i,t_{i+1}]=[x_i,x_{i+1}].
  \]

So **every** edge of \(R\) is one of the segments
\[
[x_1,x_2],\ [x_2,x_3],\ \dots,\ [x_n,x_1].
\]

Conversely, Step 2 shows each of these segments is indeed an edge of \(R\). Thus the boundary edges of \(R\) are exactly those \(n\) segments.

---

### Step 4: cyclic order of the vertices
The boundary of a convex polygon is a simple closed polygonal chain made of its edges.  
We have just proved that the boundary edges of \(R\) are exactly
\[
[x_1,x_2],\ [x_2,x_3],\ \dots,\ [x_n,x_1].
\]
Therefore the boundary of \(R\) is precisely the closed chain
\[
x_1,x_2,\dots,x_n,x_1.
\]
Hence the vertices of \(R\) are exactly
\[
x_1,x_2,\dots,x_n,
\]
they are pairwise distinct, and they occur in that cyclic order.

This proves Part \((1)\).

---

## Final conclusion

For each \(i\),
\[
x_i=s_i+t_i
\]
is a vertex of \(P+Q\), and the vertices of \(P+Q\) are exactly
\[
x_1,x_2,\dots,x_n
\]
in cyclic order.

Also, for any outer normal to the edge \([s_i,s_{i+1}]\) of \(P\),
\[
F_Q(u)=[t_i,t_{i+1}],
\]
and
\[
t_{i+1}-t_i=\lambda_i(s_{i+1}-s_i)\qquad\text{for some }\lambda_i\ge 0.
\]

So both statements are proved.